% ભાગ: first-order-logic
% પ્રકરણ: beyond
% વિભાગ: modal-logics

\documentclass[../../../include/open-logic-section]{subfiles}

\begin{document}

\olfileid{fol}{byd}{mod}

\olsection{મોડલ તર્કશાસ્ત્રો}

નીચેના શરતી વાક્યને વિચારો:
\begin{quote}
  જો જેરેમી એ ઓરડામાં એકલો હોય, તો તે દારૂ પીધેલો, નિર્વસ્ત્ર અને
  ખુરશીઓ પર નાચતો હશે.
\end{quote}
આવું શરતી વિધાન સત્યતાફલનલક્ષી અર્થમાં સત્ય હોવા છતાં ભ્રામક હોઈ શકે
છે, કારણ કે તેના શબ્દો જાણે એવું સૂચવે છે કે પૂર્વાંગ અસત્ય હોય અથવા
ફલિત સત્ય હોય એટલા કરતાં પૂર્વાંગ અને ફલિત વચ્ચે વધુ મજબૂત સંબંધ છે.
એટલે તેના શબ્દો સૂચવે છે કે આ દાવો માત્ર આ ચોક્કસ જગતમાં જ સત્ય નથી
(જ્યાં જેરેમી ઓરડામાં એકલો ન હોવાથી તે તુચ્છ રીતે સત્ય હોઈ શકે), પરંતુ
પૂર્વાંગ સત્ય \emph{હોત}, તો ફલિત પણ સત્ય \emph{હોત}. બીજા શબ્દોમાં,
આ વિધાનનો અર્થ દાવો માત્ર આ જગતમાં નહિ પણ દરેક ``શક્ય'' જગતમાં સત્ય
છે; અથવા આ ચોક્કસ જગતમાં માત્ર સત્ય હોવાને બદલે \emph{અનિવાર્યપણે}
સત્ય છે એવો લઈ શકાય.

આ પ્રકારની અનિવાર્યતાનો અર્થ સમજાવવા મોડલ તર્કશાસ્ત્ર રચાયું હતું.
સામાન્ય વિધાનાત્મક તર્કશાસ્ત્રમાં એક પેટી-કારક ઉમેરીને મોડલ
વિધાનાત્મક તર્કશાસ્ત્ર મળે છે; એટલે જો~$!A$~!!a{formula} હોય, તો
$\Box !A$ પણ સૂત્ર છે. સહજ રીતે, $\Box !A$ કહે છે કે~$!A$
\emph{અનિવાર્યપણે} સત્ય છે, અથવા દરેક શક્ય જગતમાં સત્ય છે.
$\Diamond !A$ને સામાન્ય રીતે $\lnot \Box \lnot !A$નું સંક્ષેપ માનવામાં
આવે છે અને તે~$!A$ \emph{શક્યપણે} સત્ય છે એમ કહે છે એવું વાંચી શકાય.
અલબત્ત, વિધેય તર્કશાસ્ત્રમાં પણ મોડલતા ઉમેરી શકાય છે.

મોડલ તર્કશાસ્ત્રનું અર્થવિજ્ઞાન આપવા ક્રિપ્કે !!{structure}s વાપરી
શકાય છે; હકીકતમાં ક્રિપ્કેએ આ અર્થવિજ્ઞાન સૌપ્રથમ મોડલ તર્કશાસ્ત્રને
ધ્યાનમાં રાખીને રચ્યું હતું. આંશિક ક્રમો સુધી મર્યાદિત રહેવાને બદલે
વધુ સામાન્ય રીતે ``શક્ય જગતો''નો ગણ~$P$ અને જગતો વચ્ચેનો દ્વિસ્થાની
``સુલભતા'' સંબંધ~$\Atom{R}{x,y}$ હોય છે. સહજ રીતે,
$\Atom{R}{p,q}$ કહે છે કે જગત~$q$, જગત~$p$ સાથે સુસંગત છે; એટલે જો
આપણે જગત~$p$માં ``હોઈએ'', તો જગત~$q$ જેવું પણ હોઈ શક્યું હોત એવી
શક્યતા વિચારવી પડે.

મોડલ તર્કશાસ્ત્રને કેટલીક વાર ``અંતર્લક્ષી'' તર્કશાસ્ત્ર કહે છે અને
તેનો ભેદ ``વિસ્તારલક્ષી'' તર્કશાસ્ત્રથી પાડે છે. પ્રશિષ્ટ તર્કશાસ્ત્ર
જેવા વિસ્તારલક્ષી તર્કશાસ્ત્રનું અભિપ્રેત અર્થવિજ્ઞાન માત્ર એક, એટલે
``વાસ્તવિક'', જગતનો ઉલ્લેખ કરે છે; જ્યારે ``અંતર્લક્ષી'' તર્કશાસ્ત્રનું
અર્થવિજ્ઞાન વધુ વિસ્તૃત તત્ત્વમીમાંસા પર આધાર રાખે છે. અનિવાર્યતાને
સંરચિત કરવા ઉપરાંત, ઉપયોગને અનુરૂપ $\Box$ અને~$\Diamond$નું પુનઃઅર્થઘટન
કરીને બીજા ભાષાકીય બંધારણોને સંરચિત કરવા પણ મોડલતા વાપરી શકાય છે.
દાખલા તરીકે:
\sourcecorrection{OLBYD-005}{મૂળની
\texttt{modal-logics.tex}, પંક્તિઓ~54--55માં ``structuring'' અને
``structure'' ક્રિયાપદોની અંદર શબ્દકોશ-ટોકન
\texttt{\string!\string!\{structure\}}
ઘુસાડવામાં આવ્યું છે, જે ગુજરાતી ટોકન-વિસ્તરણ સાથે વ્યાકરણવિરોધી થાય
છે. અહીં બંને ક્રિયાપદોનો અર્થ ``સંરચિત કરવું'' તરીકે સીધો અનુવાદિત
કર્યો છે.}
\begin{enumerate}
\item સાબિતીયોગ્યતા તર્કશાસ્ત્રમાં $\Box !A$ને ``$!A$ સાબિત કરી શકાય
  છે'' અને $\Diamond !A$ને ``$!A$ સુસંગત છે'' એમ વાંચવામાં આવે છે.
\item જ્ઞાનમીમાંસક તર્કશાસ્ત્રમાં $\Box !A$ને ``હું~$!A$ જાણું છું''
  અથવા ``હું~$!A$ માનું છું'' એમ વાંચી શકાય.
\item કાલિક તર્કશાસ્ત્રમાં $\Box !A$ને ``$!A$ હંમેશાં સત્ય છે'' અને
  $\Diamond !A$ને ``$!A$ કેટલીક વાર સત્ય છે'' એમ વાંચી શકાય.
\end{enumerate}

મોડલતા સાથે વ્યવહાર કરતાં નિયમો અને સ્વયંસિદ્ધિઓ વડે તર્કશાસ્ત્રને
વિસ્તારવા આપણે માગીએ. દાખલા તરીકે, તંત્ર~$\Log{S4}$માં વિધાનાત્મક
તર્કશાસ્ત્રનાં સામાન્ય સ્વયંસિદ્ધિઓ અને નિયમો સાથે નીચેની સ્વયંસિદ્ધિઓ
હોય છે:
\begin{align*}
& \Box (!A \lif !B) \lif (\Box !A \lif \Box !B)\\
& \Box !A \lif !A \\
& \Box !A \lif \Box \Box !A
\intertext{અને આ નિયમ પણ: ``$!A$માંથી $\Box !A$ તારવો.''
  $\Log{S5}$માં નીચેનું સ્વયંસિદ્ધિ ઉમેરાય છે:}
& \Diamond !A \lif \Box \Diamond !A
\end{align*}
આ સ્વયંસિદ્ધિઓનાં રૂપાંતરો જુદા ઉપયોગો માટે યોગ્ય હોઈ શકે; દાખલા તરીકે,
S5ને સામાન્ય રીતે તાર્કિક અનિવાર્યતાના ખ્યાલનું લક્ષણ માનવામાં આવે છે.
અને સારી વાત એ છે કે ક્રિપ્કે !!{structure}sના સુલભતા સંબંધને કુદરતી
રીતે મર્યાદિત કરીને સામાન્ય રીતે એવું અર્થવિજ્ઞાન મેળવી શકાય છે જેના
માટે !!{derivation} તંત્ર યથાર્થ અને પૂર્ણ હોય. દાખલા તરીકે,
$\Log{S4}$ ક્રિપ્કે !!{structure}sના એવા વર્ગને અનુરૂપ છે જેમાં સુલભતા
સંબંધ સ્વવાચક અને પરંપરિત હોય. $\Log{S5}$ ક્રિપ્કે !!{structure}sના એવા
વર્ગને અનુરૂપ છે જેમાં સુલભતા સંબંધ {\em સાર્વત્રિક} હોય, એટલે દરેક
જગત દરેક બીજા જગતમાંથી સુલભ હોય; તેથી $\Box !A$ ત્યારે અને માત્ર ત્યારે
સત્ય હોય જ્યારે~$!A$ દરેક જગતમાં સત્ય હોય.

\end{document}
